\section{Concurrency, User-Defined Linear Ghost State, and Atomics} \label{sec:concurrency}

Rust's memory safety and ownership discipline allows our verification methodology to be sound in the presence of multi-threading.
However, verifying low-level code with fine-grained concurrency still requires additional techniques.

One key such technique is \emph{user-defined ghost state}: just as \verus provides
\verb`PermData` to track memory ownership,
the user can define their own ghost state to track elements of a custom concurrent protocol.
For defining ghost state, we primarily use a technique of prior work~\cite{ironsync-tr},
which suggests viewing user-defined ghost state as a ``localized transition system.''
In a localized transition system,
the user defines state transitions that can be expressed in terms of
thread-local views of the global program state, and then proves inductive invariants
on the resulting state transition system.

The result of this construction is a collection of \verb`proof`-mode (ownership-checked)  ghost types
representing components of the system state, along with an API for performing operations that manipulate
the ghost objects (constructing them, dropping them, or modifying them).
These operations might require certain properties to hold of the ghost state,
which can be proved from the inductive invariants of the transition system.
Finally, the programmer can manipulate these objects like any other ghost object, e.g., putting them
inside cells, invariant objects (\autoref{sec:invcell}), locks, atomics, or other mechanisms.

For example, we use this technique to verify a FIFO queue using a ring buffer with
atomic head and tail pointers.
At a very high level, we do this by first defining ghost state to represent the evolution
of the FIFO state. This state includes both the head and tail pointers, and as a result,
\verus gives us access to ghost objects that represent the head and tail,
and we then associate these ghost objects with atomic memory using an \verb`AtomicInvariant`.

For example, one of the transitions defined in this system (out of four total)
is called \verb`consume_start`.
Its corresponding API function has the following type signature:
\begin{lstlisting}[language=Verus,style=VerusLineNos]
#[proof] pub fn consume_start(
    #[proof] &self, #[proof] tail: &Fifo::tail<T>, #[proof] consumer: &mut Fifo::consumer<T>
  ) -> PermData<T> { /* auto-generated by Verus ghost state machinery */ }
\end{lstlisting}
The \verb`self` object, here, is a (ghost) metadata object that gives access to the API.
The interesting parameters are the \verb`tail`, a user ghost state object that
represents the value of the tail pointer,
and \verb`consumer`, which represents the thread-local state of the consumer thread.
Intuitively, this signature requires two things:
first, that the client ``prove'' that they are the consumer
by exhibiting the ghost state thread in order to perform the action.
Second, that they access the tail pointer while performing the action.
If the value of the tail pointers indicates that a message is waiting to be received,
then the consumer thread obtains the permission to access
a cell of the ring buffer, from which it can read a message, and which it relinquishes
at the end of the ``consume'' operation.
The validity of the operation (i.e., its ability to return this particular ghost object)
is encoded in the correctness conditions of the transition system,
and \verus requires the user to prove that these conditions hold from the inductive invariants.

Though user ghost state is usually intended for concurrent code, it is sometimes useful
for single-threaded code as well.
The supplementary materials~\cite{verussupplementary} include the following collection of examples for both
single-threaded and multi-threaded code, all with user-defined ghost state:
\begin{itemize}
  \item A concurrent FIFO queue based on a ring buffer, with head and tail pointers manipulated by atomics, as described above.
  \item A string interner that returns an identifier and ghost state, allowing the user to reason about the identifier
      as if it were the originally interned value.
  \item A thread-safe reader-writer lock, also implemented with atomics, which allows the user to specify an invariant on the protected data, in a similar fashion to \verb`InvCell`.
  \item A (non-thread-safe) reference-counted pointer, similar to Rust's \verb`Rc` (though without weak-pointers),
      which uses a \verb`PPtr` for the heap allocation and a \verb`PCell` for the reference counter.
\end{itemize}
